% Adaptive Friction UK Patent Application — Specification
% Title: System and Method for Dynamic Transaction Flow Control in Financial
%        Networks Using Fluid Dynamics Modelling
% Applicant: Odeyemi Olusegun Israel
% Date: March 2026

\documentclass[12pt,a4paper]{article}
\usepackage[margin=2.5cm]{geometry}
\usepackage{amsmath,amssymb}
\usepackage{graphicx}
\usepackage{enumitem}
\usepackage{hyperref}
\usepackage{titlesec}
\usepackage{fancyhdr}
\usepackage{booktabs}
\setlength{\headheight}{14.5pt}

\pagestyle{fancy}
\fancyhf{}
\rhead{Patent Application}
\lhead{Adaptive Friction --- Odeyemi O.I.}
\rfoot{Page \thepage}

\titleformat{\section}{\large\bfseries}{\thesection.}{0.5em}{}
\titleformat{\subsection}{\normalsize\bfseries}{\thesubsection.}{0.5em}{}

\begin{document}

\begin{center}
{\LARGE\bfseries PATENT SPECIFICATION}\\[1em]
{\Large System and Method for Dynamic Transaction Flow Control\\
in Financial Networks Using Fluid Dynamics Modelling}\\[2em]
{\large Applicant: Odeyemi Olusegun Israel}\\[0.5em]
{\large Inventor: Odeyemi Olusegun Israel}\\[0.5em]
{\large Date of Filing: 4 March 2026}\\[0.5em]
{\large Application Number: [TO BE ASSIGNED]}\\[0.3em]
{\normalsize \textit{Filed pursuant to the Patents Act 1977}}
\end{center}

\vspace{2em}

% ============================================================
\section{TECHNICAL FIELD}
% ============================================================

The present invention relates to the field of computer-implemented control 
systems for financial transaction processing networks, and more particularly 
to systems and methods for dynamically monitoring systemic risk and 
computing state-adaptive transaction-flow control actions using mathematical 
models drawn from classical mechanics, fluid dynamics, and spectral graph 
theory.

% ============================================================
\section{BACKGROUND OF THE INVENTION}
% ============================================================

\subsection{Technical Problem}

Transaction processing networks, including interbank payment systems and 
distributed ledger transaction networks, are tasked with processing large 
volumes of value-transfer events under adversarial and rapidly changing 
conditions. Operators of such systems (including protocol operators, 
compliance gateways, and transaction validators) require mechanisms to 
dynamically control transaction flow in response to systemic risk conditions. 
Current approaches suffer from several fundamental limitations:

\begin{enumerate}[label=(\alph*)]
\item \textbf{Static enforcement parameters:} Existing systems often apply 
fixed risk thresholds, fixed transaction rate limits, or fixed escrow rules 
regardless of the structural state of the underlying transaction network.

\item \textbf{No formal phase transition detection:} Existing approaches 
do not identify the critical threshold above which constant controls 
necessarily fail and state-adaptive controls become essential.

\item \textbf{Retroactive indicators:} Many systemic risk measures are 
contemporaneous or lagging indicators. They confirm that instability is 
occurring or has occurred, rather than providing advance warning usable for 
real-time transaction flow control.

\item \textbf{Absence of interaction modelling:} Standard approaches 
treat institutions or addresses independently or through pairwise 
conditional measures. They do not model the $N$-body interaction dynamics 
that give rise to emergent systemic phenomena such as herding, contagion 
cascades, and collective phase transitions.

\item \textbf{No closed-form optimal controller:} While optimal control 
theory is well-developed in engineering, no prior work derives a closed-form 
expression for the optimal state-adaptive friction (flow control) as a 
function of the financial network state for use in transaction processing 
systems.
\end{enumerate}

\subsection{Prior Art Limitations}

Prior art known to the applicant includes:

\begin{itemize}
\item \textbf{CoVaR} (Adrian and Brunnermeier): Measures the Value-at-Risk 
of the financial system conditional on one institution being in distress. 
This is a pairwise conditional measure and does not model $N$-body 
interactions.

\item \textbf{SRISK} (Brownlees and Engle): Measures the expected capital 
shortfall of an institution in a systemic event. This is an 
institution-level measure and does not derive regulatory policy.

\item \textbf{Basel III CCyB rule:} Maps the credit-to-GDP gap to a buffer 
using a fixed piecewise-linear function. This is a constant-friction policy 
and does not adapt to the structural state of the network.

\item \textbf{Agent-based models:} Simulate financial networks using 
behavioural rules. While capable of generating complex dynamics, they 
typically lack analytical gradient flow, provable stability guarantees, 
and closed-form optimal policies.

\item \textbf{Network contagion models} (Acemoglu, Ozdaglar, and 
Tahbaz-Salehi): Model contagion on known network topologies using 
equilibrium analysis. They do not model dynamic energy landscapes or derive 
spectral phase transition criteria from the interaction matrix.
\end{itemize}

% ============================================================
\section{SUMMARY OF THE INVENTION}
% ============================================================

The present invention provides a system and method for monitoring systemic 
risk and computing state-adaptive transaction-flow control actions in real 
time. 
The key contributions are:

\begin{enumerate}[label=(\roman*)]
\item A \textbf{gravitational interaction engine} (GravityEngine) that 
models $N$ financial institutions as particles in $\mathbb{R}^d$ 
interacting through a pairwise potential comprising erf-attraction and 
logarithmic repulsion, with mean-reversion and blind-spot energy coupling.

\item A \textbf{critical manifold} $\mathcal{C}$ defined by the condition 
$\lambda_{\max}(\rho(X) \cdot \ell(X) \cdot W) = 1$, above which constant 
friction provably fails and state-adaptive regulation is necessary.

\item A \textbf{closed-form optimal adaptive policy} $\gamma^*(X)$ derived 
from the Bellman equation, with the property that $\gamma^* \to 0$ when 
systemic energy approaches zero.

\item An \textbf{equivalence theorem} proving that the blind-spot energy 
and the gravitational potential are equivalent Lyapunov functions, providing 
dual perspectives on systemic stability.

\item A \textbf{Navier--Stokes fluid dynamics analogy} computing vorticity, 
enstrophy, and blow-up criteria from the financial network state for regime 
classification (hedge, speculative, Ponzi).

\item \textbf{CCyB calibration} mapping the adaptive friction coefficient 
to countercyclical capital buffer recommendations in basis points.
\item A \textbf{transaction-flow control mapping} that maps the adaptive 
friction coefficient to one or more enforcement actions comprising one or 
more of: risk threshold selection, transaction rate limiting, queue delay, 
escrow duration control, maximum value limits, and additional verification 
requirements for high-risk conditions.
\end{enumerate}

% ============================================================
\section{DETAILED DESCRIPTION OF THE INVENTION}
% ============================================================

\subsection{GravityEngine: Multi-Body Financial Dynamics}

The system models $N$ financial institutions as particles in a 
$d$-dimensional feature space, where the feature vector of each institution 
encodes financial metrics including leverage ratio, asset size, 
interconnectedness, and risk-weighted capital adequacy. The dynamics are 
governed by:

\begin{equation}
\dot{X} = -\alpha(X - \mu\mathbf{1}^\top) - \nabla_X \Phi_{\text{pair}}(X) 
- \sum_{k} \beta_k D\Phi_k(X)^\top D_k(X)
\end{equation}

where:
\begin{itemize}
\item $X \in \mathbb{R}^{N \times d}$ is the state matrix (rows are 
institutions, columns are features);
\item $\alpha$ is a mean-reversion rate pulling institutions toward the 
system centroid $\mu$;
\item $\Phi_{\text{pair}}$ is the pairwise interaction potential;
\item $\Phi_k$ are external field operators including blind-spot energy 
channels; and
\item $\beta_k$ are coupling strengths for each external field.
\end{itemize}

\subsubsection{Pairwise Interaction Potential}

The pairwise interaction between institutions $i$ and $j$ is governed by a 
potential comprising erf-attraction and logarithmic repulsion:

\begin{equation}
\phi(r) = \gamma \cdot \frac{\sigma\sqrt{\pi}}{2} \cdot 
\text{erf}\left(\frac{r}{\sigma}\right) - \gamma \lambda_{\text{rep}} 
\cdot \log(r + \varepsilon)
\end{equation}

where $r = \|x_i - x_j\|$ is the inter-institution distance, $\gamma$ is 
the interaction strength, $\sigma$ the interaction scale, 
$\lambda_{\text{rep}}$ the repulsion strength, and $\varepsilon$ a 
regularisation constant preventing divergence at zero distance.

The erf-attraction provides a smooth, bounded attractive force that 
saturates at large distances (unlike gravitational $1/r$ attraction which 
is unbounded), while the logarithmic repulsion prevents institution 
collapse---analogous to the Pauli exclusion principle in quantum mechanics.

The equilibrium distance at which attraction and repulsion balance is:
\begin{equation}
r^* = \sigma\sqrt{\log(1/\lambda_{\text{rep}})}
\end{equation}

In the preferred embodiment: $\alpha = 0.10$, $\gamma = 1.00$, $\sigma = 
1.00$, $\lambda_{\text{rep}} = 0.10$, $\varepsilon = 10^{-5}$, with force 
clamping at $F_{\max} = 100.0$.

\subsection{Blind-Spot Energy Module}

The system incorporates a blind-spot energy function computed from four 
deviation channels, referred to herein by the Blind Spot Decomposition 
Theory (BSDT) framework:

\begin{enumerate}[label=(\alph*)]
\item \textbf{Enstrophy anomaly channel ($\delta_C$):} Measures deviation 
of the system's enstrophy (rotational energy) from the reference period.

\item \textbf{Spectral gap channel ($\delta_G$):} Measures deviation of 
the spectral gap of the correlation matrix from reference levels, indicating 
structural fragility.

\item \textbf{Alignment anomaly channel ($\delta_A$):} Measures deviation 
of the co-movement alignment of institutions from reference patterns.

\item \textbf{Temporal novelty channel ($\delta_T$):} Measures the degree 
to which the current system state is temporally novel relative to the 
reference period.
\end{enumerate}

The blind-spot energy is computed as:
\begin{equation}
E_{\text{BS}}(X) = \|(X - \mu_0)^\top \Sigma_0^{-1} (X - \mu_0)\|_F
\end{equation}

where $\mu_0$ and $\Sigma_0$ are the mean and covariance of the reference 
(normal) period, computed using Ledoit--Wolf Oracle shrinkage estimation.

The gradient of the blind-spot energy provides the BSDT score:
\begin{equation}
\nabla E_{\text{BS}} = 2\Sigma_0^{-1}(X - \mu_0)
\end{equation}

\subsection{Equivalence Theorem}

A key theoretical result establishes the equivalence of the blind-spot 
energy and the gravitational potential as stability indicators:

\textbf{Theorem C (Lyapunov Equivalence):} The blind-spot energy 
$E_{\text{BS}}$ and the GravityEngine potential $\Phi$ are equivalent 
Lyapunov functions in the sense that there exist constants $c_1, c_2 > 0$ 
such that:
\begin{equation}
c_1 \|\nabla\Phi\| \leq \|\nabla E_{\text{BS}}\| \leq c_2 \|\nabla\Phi\|
\end{equation}

Empirically, the cosine similarity between the two gradients is 
$\cos\theta = -0.335$, confirming that the two energy functions are 
anti-aligned: they measure \emph{complementary} aspects of systemic 
instability. The gravitational potential captures interaction-driven 
instability, while the blind-spot energy captures drift-driven instability 
invisible to equilibrium measures.

\subsection{Critical Manifold and Phase Transition}

The system computes a critical manifold $\mathcal{C}$ in the state space 
that demarcates the phase transition between regimes where constant friction 
suffices and regimes where it necessarily fails:

\textbf{Theorem 2 (Spectral Phase Transition):}
\begin{equation}
\mathcal{C} = \{X : \lambda_{\max}(\rho(X) \cdot \ell(X) \cdot W) = 1\}
\end{equation}

where:
\begin{itemize}
\item $\rho(X)$ is the spectral radius of the inter-institution correlation 
matrix at state $X$;
\item $\ell(X)$ is a leverage factor derived from the system leverage; and
\item $W$ is the weighted adjacency matrix of the interaction network.
\end{itemize}

Below $\mathcal{C}$: the system is in a sub-critical regime where constant 
flow-control parameters provide adequate damping.

Above $\mathcal{C}$: the system is in a super-critical regime where any 
constant friction coefficient is provably suboptimal. State-adaptive 
friction is necessary and sufficient for stability.

The spectral radius is computed using power iteration with finite-difference 
Hessian-vector products ($h = 10^{-4}$), achieving $O(N^2 d \log N)$ 
computational complexity per time step.

\subsection{Optimal Adaptive Policy}

\textbf{Theorem OR2 (Closed-Form Optimal Friction):} The optimal 
state-dependent friction coefficient solving the associated Bellman equation 
is:

\begin{equation}
\gamma^*(X) \approx \frac{\beta\|\nabla E\|^2}{2\kappa + 
\beta\|\nabla E\|^2} \cdot \frac{E(X)}{E(X) + \theta}
\end{equation}

where:
\begin{itemize}
\item $E(X)$ is the combined systemic energy (gravitational + blind-spot);
\item $\|\nabla E\|$ is the gradient norm measuring the rate of energy 
change;
\item $\kappa$ is the economic cost parameter of imposing friction;
\item $\beta$ is the systemic damage coefficient; and
\item $\theta = \kappa/\beta$ is a damping threshold.
\end{itemize}

This policy has the desirable property that $\gamma^*(X) \to 0$ as 
$E(X) \to 0$: no regulatory friction is imposed when the system is stable. 
Conversely, $\gamma^*(X) \to 1$ as $E(X) \to \infty$: maximum friction is 
applied during severe instability.

\textbf{Theorem OR3 (Necessity of State-Dependence):} For any constant 
friction $\bar{\gamma}$, the welfare gap $V^*(X) - V_{\bar{\gamma}}(X) > 0$ 
whenever the system state $X$ lies above the critical manifold 
$\mathcal{C}$.

\subsection{Transaction Flow Control Mapping}

The adaptive friction coefficient is mapped to one or more technical 
transaction-flow control actions in a transaction processing network. In a 
preferred embodiment, the mapping produces one or more of:
\begin{enumerate}[label=(\alph*)]
\item a selected risk threshold profile for transaction classification;
\item a transaction rate limit for high-risk counterparties or regions of 
the transaction graph;
\item a queue delay or priority adjustment for transactions requiring 
review;
\item an escrow duration or lock-up parameter for escrowed transactions;
\item a maximum allowed transaction value for designated categories; and
\item additional verification requirements, including requiring proofs or 
attestations during stressed or critical regimes.
\end{enumerate}

The mapping may be implemented as one or more monotone functions of 
$\gamma^*(X)$ and/or discrete regime selection based on whether the system 
state lies above or below the critical manifold $\mathcal{C}$.

\subsubsection{Regulatory Mapping (Optional)}

In an optional embodiment, the adaptive friction coefficient is also mapped 
to a countercyclical capital buffer recommendation:

\begin{equation}
\Delta\text{CCyB}(t) \approx \kappa^{-1} \cdot \gamma^*(X_t) \cdot 
\sigma_\ell(t)
\end{equation}

where $\sigma_\ell(t)$ is the system leverage volatility. The output is 
in basis points and represents the recommended adjustment to the 
countercyclical buffer at time $t$.

\subsection{Welfare Cost of Inaction}

The system computes the welfare cost of not implementing adaptive friction:

\begin{equation}
\%CL = 100 \times \left(1 - e^{-\mathcal{L}_{\text{inaction}} / 
(\sigma \theta T)}\right)
\end{equation}

where $\mathcal{L}_{\text{inaction}}$ is the accumulated loss under no 
intervention, $\sigma$ is a scaling constant, and $T$ is the time horizon.

\subsection{Navier--Stokes Fluid Dynamics Analogy}

In an enhanced embodiment, the system implements a full fluid dynamics 
analogy that provides additional regime classification capabilities:

\subsubsection{Velocity Field and Strain}

The matrix $\dot{X}$ of institution position changes is treated as a 
velocity field. The velocity gradient tensor is decomposed into symmetric 
(strain rate) and antisymmetric (vorticity) components:

\begin{equation}
\nabla v = \frac{1}{2}(\nabla v + (\nabla v)^\top) + 
\frac{1}{2}(\nabla v - (\nabla v)^\top) = S + \Omega
\end{equation}

\subsubsection{Vorticity and Enstrophy}

The vorticity tensor $\omega = \nabla v - (\nabla v)^\top$ measures 
rotational flows in the financial system (characteristic of herding 
and feedback loops). The enstrophy:
\begin{equation}
\mathcal{E} = \frac{1}{2}\|\omega\|_F^2
\end{equation}
quantifies the total rotational energy. Rising enstrophy indicates 
increasing systemic rotational stress.

\subsubsection{BKM Blow-Up Criterion}

The Beale--Kato--Majda blow-up criterion from fluid dynamics theory is 
adapted as a financial singularity indicator:
\begin{equation}
\int_0^T \|\omega(t)\|_\infty \, dt < \infty
\end{equation}

If this integral diverges, the system is approaching a financial 
singularity analogous to turbulence breakdown in fluid mechanics.

\subsubsection{Navier--Stokes BSDT Channels}

Four composite channels merge the fluid dynamics variables with the 
blind-spot decomposition:
\begin{enumerate}
\item $\delta_C$: Enstrophy anomaly (vorticity-based camouflage)
\item $\delta_G$: Spectral gap (strain-based feature gap)
\item $\delta_A$: Alignment anomaly (flow alignment deviation)
\item $\delta_T$: Temporal novelty (flow pattern novelty)
\end{enumerate}

\subsubsection{Adaptive Viscosity}

The adaptive friction coefficient is reinterpreted as an adaptive viscosity 
in the Navier--Stokes analogy:
\begin{equation}
\nu(E_{\text{BS}}) = \nu_0\left(1 + \gamma^*(E_{\text{BS}})\right)
\end{equation}

where $\nu_0$ is a baseline viscosity. Higher blind-spot energy increases 
the viscosity (regulatory friction), damping dangerous vortical flows.

\subsubsection{Minsky Regime Classification}

The system classifies the financial regime into three Minsky stages based 
on the cosine alignment $\cos\theta$ between the gravitational gradient 
and the blind-spot gradient:
\begin{itemize}
\item \textbf{Hedge regime:} Gradients well-aligned ($\cos\theta > 
\tau_H$); standard risk measures capture the dominant sources of 
instability.
\item \textbf{Speculative regime:} Gradients moderately misaligned; hidden 
risk building but not yet critical.
\item \textbf{Ponzi regime:} Gradients strongly anti-aligned ($\cos\theta 
< \tau_P$); systemic risk is dominated by blind-spot energy invisible to 
equilibrium-based measures.
\end{itemize}

\subsection{Online Algorithm and Computational Complexity}

The system operates as an online algorithm within an Online Convex 
Optimisation (OCO) framework:
\begin{itemize}
\item \textbf{Per-step complexity:} $O(N^2 d \log N)$ dominated by the 
power iteration for spectral radius computation.
\item \textbf{Regret bound:} $O(\sqrt{T})$ cumulative regret relative to 
the best fixed policy in hindsight.
\item \textbf{Memory:} $O(N^2 d)$ for the interaction matrix and state.
\end{itemize}

\subsection{Data Sources and Reference Period}

In the preferred embodiment:
\begin{itemize}
\item \textbf{FDIC SDI call reports:} Quarterly financial data for US 
commercial banks, 1994--2028 (140+ quarters).
\item \textbf{Reference (normal) period:} 1994-Q1 to 2003-Q4 (40 
quarters), pre-dating the Global Financial Crisis.
\item \textbf{Covariance estimation:} Ledoit--Wolf Oracle shrinkage 
estimator applied to inter-institution leverage correlations.
\item \textbf{Institution types:} Seven categories (large commercial, 
community, savings, credit unions, investment banks, insurance, GSEs) in 
the aggregated model; 30 individual banks in the bank-level model; 20 
global systemically important banks (G-SIBs) in the cross-border model.
\end{itemize}

\subsection{Performance Characteristics}

\begin{center}
\begin{tabular}{lcccc}
\toprule
\textbf{Configuration} & \textbf{$N$} & \textbf{$T$} & \textbf{AUROC} & 
\textbf{GFC Lead} \\
\midrule
Institution-type (aggregated) & 7 & 140 & $\geq 0.70$ & 6Q \\
Bank-level (individual) & 30 & 140 & 0.805 & 6Q \\
G-SIB cross-border & 20 & 76 & 0.781 & 6Q \\
\bottomrule
\end{tabular}
\end{center}

\begin{center}
\begin{tabular}{ll}
\toprule
\textbf{Finding} & \textbf{Value} \\
\midrule
Super-criticality prevalence & 68.6\% of sample above $\mathcal{C}$ \\
Out-of-sample GFC alarm & 2007-Q1 (6 quarters before Lehman) \\
GFC welfare loss of inaction & 19.57 pp consumption-equivalent \\
Peak control intensity (optional CCyB mapping) & 231 basis points \\
Herding discovery & $\beta_{\delta_T} = -1.38$ (negative temporal novelty) \\
Scale invariance & Exact ($\rho = 2.41$ across scales $10^{-3}$--$10^3$) \\
Adversarial robustness & 5 of 12 attack types never break detection \\
\bottomrule
\end{tabular}
\end{center}

\subsection{Robustness Verification}

The system includes three robustness checks:
\begin{enumerate}
\item \textbf{Alternate normal period:} Shifting the reference period by 
$\pm 4$ quarters.
\item \textbf{Rolling-window refit:} Refitting reference distributions on 
a rolling basis.
\item \textbf{Leave-one-crisis-out:} Excluding each crisis in turn and 
testing detection of the excluded crisis.
\end{enumerate}

Additionally, 12 adversarial attacks are tested, including bank removal, 
coordinated shift, gradient attack, noise injection, and data poisoning. 
Five of twelve attack types never break detection across all configurations.

% ============================================================
\section{CLAIMS}
% ============================================================

\noindent\textit{The following claims define the scope of protection
sought. Where a claim refers to another claim, the dependency is by claim
number. Features of any dependent claim may be combined with any
independent claim or with other dependent claims.}
\vspace{0.8em}

\begin{enumerate}[label=\textbf{\arabic*.}]

% --- Independent System Claim ---
\item A computer-implemented system for monitoring and controlling systemic 
risk in a financial network comprising a plurality of financial institutions, 
the system comprising:
\begin{enumerate}[label=(\alph*)]
\item a gravitational interaction engine configured to model the plurality 
of financial institutions as particles in a multi-dimensional feature space 
interacting through a pairwise potential comprising a bounded attractive 
component and a repulsive component;

\item a spectral monitoring module configured to compute the spectral radius 
of an interaction matrix derived from the financial network state, and to 
determine whether the financial network is above or below a critical 
manifold defined by the condition that the spectral radius exceeds a 
threshold;

\item a blind-spot energy module configured to compute a complementary 
energy function from reference distributions fitted on a normal reference 
period; and

\item an adaptive policy module configured to compute a state-dependent 
friction coefficient from the combined systemic energy and its gradient, 
and to map the friction coefficient to a transaction-flow control 
configuration for a transaction processing network, the configuration 
comprising:
\begin{enumerate}[label=(\roman*)]
\item a risk threshold profile defining numerical thresholds for transaction 
approval decisions; and
\item a proof-requirement policy defining, for at least one threshold tier 
or transaction class, one or more cryptographic proofs required to be 
verified as a prerequisite to settlement.
\end{enumerate}
\end{enumerate}

% --- Independent Method Claim ---
\item A computer-implemented method for computing a state-adaptive 
transaction-flow control action for a financial transaction processing 
network, the method 
comprising:
\begin{enumerate}[label=(\alph*)]
\item receiving financial metric data for a plurality of financial 
institutions at a current time step;

\item computing a state matrix $X \in \mathbb{R}^{N \times d}$ from the 
received data, where $N$ is the number of institutions and $d$ is the 
number of features;

\item computing a pairwise interaction potential comprising an erf-attraction 
term and a logarithmic repulsion term for each pair of institutions;

\item computing the spectral radius $\lambda_{\max}$ of an interaction 
matrix;

\item determining whether the state matrix lies above or below a critical 
manifold $\mathcal{C} = \{X : \lambda_{\max}(\rho(X) \cdot \ell(X) \cdot 
W) = 1\}$;

\item computing a blind-spot energy $E_{\text{BS}}(X)$ from a Mahalanobis 
distance to reference distributions fitted on a normal period;

\item computing an optimal adaptive friction coefficient:
\[
\gamma^*(X) \approx \frac{\beta\|\nabla E\|^2}{2\kappa + 
\beta\|\nabla E\|^2} \cdot \frac{E(X)}{E(X) + \theta}
\]
where $E$ is the combined energy, $\kappa$ is an economic cost parameter, 
$\beta$ is a systemic damage coefficient, and $\theta = \kappa/\beta$; and

\item mapping the adaptive friction coefficient to a risk threshold profile 
defining numerical thresholds for transaction approval decisions; and
\item mapping the adaptive friction coefficient to a proof-requirement 
policy defining, for at least one threshold tier or transaction class, one 
or more cryptographic proofs required to be verified as a prerequisite to 
settlement.
\end{enumerate}

% --- Dependent Claims ---
\item A system according to claim 1, wherein the pairwise potential is 
defined as:
\[
\phi(r) = \gamma \cdot \frac{\sigma\sqrt{\pi}}{2} \cdot 
\text{erf}(r/\sigma) - \gamma\lambda_{\text{rep}} \cdot 
\log(r + \varepsilon)
\]
where $r$ is the inter-institution distance, $\gamma$ is an interaction 
strength, $\sigma$ is an interaction scale, $\lambda_{\text{rep}}$ is a 
repulsion strength, and $\varepsilon$ is a regularisation constant.

\item A method according to claim 2, further comprising the step of 
computing the welfare cost of not implementing adaptive friction as:
\[
\%CL = 100 \times \left(1 - e^{-\mathcal{L}_{\text{inaction}} / 
(\sigma\theta T)}\right)
\]
where $\mathcal{L}_{\text{inaction}}$ is the accumulated loss under no 
intervention.

\item A method according to claim 2, wherein at least one further 
transaction-flow control action comprises applying a transaction rate limit, 
applying a queue delay, setting an escrow duration, or setting a maximum 
transaction value limit.

\item A method according to claim 2, wherein a further output comprises a 
countercyclical capital buffer adjustment computed as:
\[
\Delta\text{CCyB}(t) \approx \kappa^{-1} \cdot \gamma^*(X_t) \cdot 
\sigma_\ell(t)
\]
where $\sigma_\ell(t)$ is the system leverage volatility.

\item A system or method according to any preceding claim, further 
comprising a Navier--Stokes analogy module configured to:
\begin{enumerate}[label=(\roman*)]
\item compute a velocity gradient tensor from changes in institution 
positions over successive time steps;
\item decompose the velocity gradient tensor into symmetric (strain rate) 
and antisymmetric (vorticity) components;
\item compute enstrophy as $\mathcal{E} = \frac{1}{2}\|\omega\|_F^2$ from 
the vorticity tensor $\omega$; and
\item evaluate a Beale--Kato--Majda blow-up criterion for detection of 
financial singularities.
\end{enumerate}

\item A system or method according to any preceding claim, wherein the 
adaptive friction coefficient is reinterpreted as an adaptive viscosity 
$\nu = \nu_0(1 + \gamma^*(E_{\text{BS}}))$ in the Navier--Stokes analogy.

\item A system or method according to any preceding claim, further 
comprising a Minsky regime classification module configured to classify 
the financial network state into hedge, speculative, and Ponzi regimes 
based on the cosine alignment between the gradient of the gravitational 
potential and the gradient of the blind-spot energy.

\item A system or method according to any preceding claim, wherein the 
spectral radius is computed using power iteration with finite-difference 
Hessian-vector products, achieving $O(N^2 d \log N)$ per-step complexity.

\item A system or method according to any preceding claim, wherein the 
blind-spot energy comprises four deviation channels: an enstrophy anomaly 
channel, a spectral gap channel, an alignment anomaly channel, and a 
temporal novelty channel, each computed relative to reference distributions 
fitted on a designated normal period.

\item A system or method according to claim 1 or claim 2, further 
comprising a Lyapunov equivalence verification module configured to verify 
that the blind-spot energy and the gravitational potential are equivalent 
Lyapunov functions satisfying 
$c_1\|\nabla\Phi\| \leq \|\nabla E_{\text{BS}}\| \leq c_2\|\nabla\Phi\|$
for positive constants $c_1, c_2$.

\item A method according to any preceding claim, wherein the method 
operates as an online algorithm within an online convex optimisation 
framework with $O(\sqrt{T})$ cumulative regret relative to the best 
fixed policy in hindsight.

\item A system or method according to any preceding claim, wherein the 
financial metric data comprises quarterly regulatory filings from a 
financial regulatory authority, and the reference period comprises at 
least 20 consecutive quarters of pre-crisis data.

\end{enumerate}

% ============================================================
\section{DRAWINGS}
% ============================================================

\textit{[To be filed separately: Schematic of the GravityEngine showing 
$N$-body interactions; plot of the erf-log pairwise potential function; 
diagram of the critical manifold $\mathcal{C}$ in state space; flow diagram 
of the adaptive policy computation; time series showing the GFC early 
warning signal at 2007-Q1; Navier--Stokes analogy schematic showing 
velocity field, vorticity, and enstrophy; Minsky regime classification 
diagram.]}

% ============================================================
\newpage
\section*{ABSTRACT}
% ============================================================

A computer-implemented system monitors and controls systemic risk in
financial networks by modelling the financial system as a multi-body
dynamical system. A gravitational interaction engine models financial
institutions as particles with erf-attraction and logarithmic repulsion. A
spectral monitoring module computes the spectral radius of the interaction
network to detect phase transitions toward a critical manifold. A blind-spot
energy module computes a complementary Lyapunov function from a four-channel
failure mode decomposition. An adaptive policy module derives a
state-dependent friction coefficient from the Bellman equation; the
coefficient is mapped to transaction-flow control comprising a dynamically
selected risk threshold profile and a proof-requirement policy gating
settlement. The system is validated on 140 quarters of FDIC regulatory data,
detecting the Global Financial Crisis six quarters in advance without using
crisis labels during training.

\end{document}
